% @@CORSO@@: appunti ad alta leggibilità per studenti con DSA
% Compilare con LuaLaTeX: latexmk -lualatex main.tex
\documentclass[12pt,a4paper,openany]{book}

\usepackage[a4paper,margin=2.5cm,headheight=16pt]{geometry}

% ======================= PROFILO DI LETTURA ==========================================
% Qui si regola tutta la leggibilità: i DSA sono diversi tra loro, quindi conviene provare.
\newcommand{\interlinea}{1.4}      % distanza tra le righe (1.2 più compatto, 1.6 più arioso)
\newcommand{\spaziolettere}{1.5}   % spazio tra le lettere, in millesimi circa (0 per toglierlo)
\newcommand{\spazioparole}{1.15}   % spazio tra le parole (1 è il valore normale)
\newcommand{\fontprincipale}{Fira Sans} % alternative: Atkinson Hyperlegible Next, Source Sans 3, Lato
\newcommand{\fontcodice}{Fira Mono}     % alternativa: Atkinson Hyperlegible Mono
\newif\ifsfondocrema
\sfondocrematrue                   % \sfondocremafalse per lo sfondo bianco, utile in stampa
% =====================================================================================

% --- Lingue ---------------------------------------------------------------------
% Aggiungi le lingue del materiale prima di main=, per esempio:
%   \usepackage[ngerman,latin,french,spanish,english,main=italian]{babel}
% Il greco, anche politonico, si scrive direttamente con Fira Sans: non serve dichiararlo.
\usepackage[english,main=italian]{babel}
% Niente legature: fi, fl e simili restano lettere separate, a favore di lettura e sintesi vocale.
\babelfont{rm}[Ligatures=NoCommon,LetterSpace=\spaziolettere,WordSpace=\spazioparole]{\fontprincipale}
\babelfont{sf}[Ligatures=NoCommon,LetterSpace=\spaziolettere,WordSpace=\spazioparole]{\fontprincipale}
\babelfont{tt}{\fontcodice}
% Per l'ebraico aggiungi bidi=basic alle opzioni di babel e togli il commento a queste righe:
% \babelprovide[import]{hebrew}
% \babelfont[hebrew]{rm}[Renderer=HarfBuzz]{Libertinus Serif}

% --- Pacchetti --------------------------------------------------------------------
\usepackage{amsmath}
\usepackage{unicode-math}
\setmathfont{Fira Math}
% Frazioni e indici a grandezza piena anche dentro il testo: aiuta in caso di discalculia.
\everymath{\displaystyle}
\usepackage{setspace}
\setstretch{\interlinea}
\usepackage{ragged2e}
\usepackage{graphicx}
\graphicspath{{images/}}
\usepackage{float}
\usepackage{etoolbox}
\usepackage[font=normalsize,labelfont=bf,justification=RaggedRight,singlelinecheck=false]{caption}
\usepackage{booktabs,tabularx,array}
\usepackage{enumitem}
\setlist{itemsep=4pt,topsep=4pt}
\usepackage{paracol}
\usepackage{xcolor}
\usepackage[most]{tcolorbox}
\usepackage{listings}
\usepackage{tikz}
\usetikzlibrary{arrows.meta,positioning,shapes.geometric,shapes.misc,automata,babel}
\usepackage{imakeidx}
\makeindex[name=indice,title=Indice analitico,intoc,columns=1]
\usepackage{fancyhdr}
\usepackage[hidelinks,unicode,bookmarksnumbered,bookmarksopen,pdflang={it-IT},pdftitle={@@CORSO@@: appunti}]{hyperref}

% --- Colori: contrasto alto ma non massimo. Il colore non è mai l'unico segnale. --------
\ifsfondocrema
  \definecolor{pagina}{HTML}{FBF7EC}
\else
  \definecolor{pagina}{HTML}{FFFFFF}
\fi
\definecolor{testo}{HTML}{222222}
\definecolor{accento}{HTML}{1F4E79}
\definecolor{corr}{HTML}{A4262C}
\definecolor{libro}{HTML}{1F6F5C}
\definecolor{fonte}{HTML}{7A5C00}
\definecolor{sintesi}{HTML}{4A4A6A}
\definecolor{mappacol}{HTML}{0B5C7A}
\definecolor{flussocol}{HTML}{6A3D7C}
\definecolor{ricorda}{HTML}{B3541E}
\definecolor{semplice}{HTML}{2F6B2F}
\pagecolor{pagina}

% --- Testo: allineato a sinistra, senza sillabazione, senza rientri, con paragrafi staccati ---
\raggedbottom
\setlength{\parindent}{0pt}
\setlength{\parskip}{0.7em plus 0.1em}
\setlength{\RaggedRightParindent}{0pt}
\setlength{\RaggedRightParfillskip}{0pt plus 1fil}
\AtBeginDocument{\color{testo}\RaggedRight\hyphenpenalty=10000\exhyphenpenalty=10000\sloppy}
% I titoli di opere e le parole straniere si distinguono con il colore, non con il corsivo.
\renewcommand{\emph}[1]{{\color{accento}#1}}
% Colonna di tabella che va a capo a sinistra, senza sillabare.
\newcolumntype{Y}{>{\RaggedRight\arraybackslash}X}
\renewcommand{\arraystretch}{1.3}

\pagestyle{fancy}
\fancyhf{}
\fancyhead[LE,RO]{\thepage}
\fancyhead[RE]{\ultimocapitolo}
\fancyhead[LO]{\ultimasezione}
\renewcommand{\headrulewidth}{0.3pt}
% Le intestazioni leggono queste macro e non \leftmark e \rightmark, che con paracol
% restano vuoti o arretrati di un capitolo.
\makeatletter
\newcommand{\ultimocapitolo}{}
\newcommand{\ultimasezione}{}
\renewcommand{\chaptermark}[1]{%
  \if@mainmatter\xdef\ultimocapitolo{\@chapapp\ \thechapter.\noexpand\ \unexpanded{#1}}%
  \else\gdef\ultimocapitolo{#1}\fi
  \gdef\ultimasezione{}\markboth{\ultimocapitolo}{\ultimasezione}}
\renewcommand{\sectionmark}[1]{%
  \xdef\ultimasezione{\thesection.\noexpand\ \unexpanded{#1}}\markright{\ultimasezione}}
\makeatother

% Fira Mono non ha il corsivo: il codice si distingue con il colore.
\lstset{basicstyle=\ttfamily,breaklines=true,columns=fullflexible,keepspaces=true,
  frame=leftline,framerule=1.5pt,rulecolor=\color{accento},xleftmargin=8pt,
  numberstyle=\small\color{gray},captionpos=b,
  keywordstyle=\color{accento},commentstyle=\color{gray},stringstyle=\color{libro}}

% --- Riquadri ---------------------------------------------------------------------------
% Ogni riquadro ha un'etichetta scritta, così si riconosce anche senza colori.
\tcbset{appbox/.style={enhanced,breakable,boxrule=0pt,leftrule=3pt,sharp corners,
  left=9pt,right=9pt,top=6pt,bottom=6pt,before skip=11pt,after skip=11pt,
  fonttitle=\bfseries,coltitle=black,detach title,halign=flush left,
  before upper={\setlength{\parskip}{.5em}\color{testo}{\tcbtitle\par}}}}
% Provenienza delle informazioni
\newtcolorbox{correzione}[1][]{appbox,colframe=corr,colback=corr!6,title={\textcolor{corr}{Correzione}},#1}
\newtcolorbox{dubbio}[1][]{appbox,colframe=fonte,colback=fonte!7,title={\textcolor{fonte}{Da verificare}},#1}
\newtcolorbox{dallibro}[1][]{appbox,colframe=libro,colback=libro!7,title={\textcolor{libro}{Dal libro}},#1}
% Struttura fissa del capitolo
\newtcolorbox{primadiiniziare}[1][]{appbox,colframe=accento,colback=accento!6,title={\textcolor{accento}{Prima di iniziare}},#1}
\newtcolorbox{mappa}[1][]{appbox,unbreakable,colframe=mappacol,colback=mappacol!4,title={\textcolor{mappacol}{Mappa del capitolo}},#1}
\newtcolorbox{daricordare}[1][]{appbox,colframe=ricorda,colback=ricorda!7,title={\textcolor{ricorda}{Da ricordare}},#1}
\newtcolorbox{insintesi}[1][]{appbox,colframe=sintesi,colback=sintesi!7,title={\textcolor{sintesi}{In sintesi}},#1}
\newtcolorbox{ripasso}[1][]{appbox,colframe=accento,colback=accento!4,title={\textcolor{accento}{Ripasso: prova a rispondere senza guardare}},#1}
\newtcolorbox{risposte}[1][]{appbox,colframe=semplice,colback=semplice!5,title={\textcolor{semplice}{Risposte, con la sezione in cui ritrovarle}},#1}
% Schemi: \begin{schema}[titolo] e \begin{flusso}[titolo]. Non si spezzano tra due pagine.
\newtcolorbox{schema}[1][]{appbox,unbreakable,colframe=mappacol,colback=mappacol!4,title={\textcolor{mappacol}{Schema: #1}}}
\newtcolorbox{flusso}[1][]{appbox,unbreakable,colframe=flussocol,colback=flussocol!4,title={\textcolor{flussocol}{Diagramma di flusso: #1}}}
% Aiuti alla comprensione: si aggiungono al testo e non lo sostituiscono.
\newtcolorbox{parolesemplici}[1][]{appbox,colframe=semplice,colback=semplice!6,title={\textcolor{semplice}{In parole semplici (spiegazione di questi appunti)}},#1}
\newtcolorbox{nonconfondere}[1][]{appbox,colframe=corr,colback=pagina,title={\textcolor{corr}{Attenzione a non confondere}},#1}
\newtcolorbox{passoapasso}[1]{appbox,colframe=flussocol,colback=flussocol!4,title={\textcolor{flussocol}{Passo per passo: #1}}}
\newtcolorbox{permemoria}[1][]{appbox,colframe=sintesi,colback=pagina,title={\textcolor{sintesi}{Per ricordarlo (trucco di questi appunti, non del corso)}},#1}

% Sezione senza slide di riferimento.
\newcommand{\soloappunti}{\par\textcolor{fonte}{\bfseries Fonte: solo appunti in classe (nessuna slide)}\par}
% Aggiunta di conoscenza generale, non dal materiale del corso.
\newcommand{\integrazione}[1]{\textcolor{sintesi}{[\textbf{Integrazione, non dal materiale del corso:} #1]}}
% Pronuncia di un nome difficile, da usare solo se è certa.
\newcommand{\pron}[1]{\textcolor{sintesi}{(si pronuncia: #1)}}

% --- Tabelle pronte --------------------------------------------------------------------
% \begin{simboli} $v$ & velocità & m/s \\ \end{simboli}: spiegazione dei simboli di una formula
\newenvironment{simboli}{\par\tabularx{\linewidth}{@{}lYl@{}}\toprule\textbf{Simbolo} & \textbf{Che cosa indica} & \textbf{Unità} \\ \midrule}{\bottomrule\endtabularx\par}
% \begin{cronologia} 1813 & evento \\ \end{cronologia}
\newenvironment{cronologia}{\par\tabularx{\linewidth}{@{}lY@{}}\toprule\textbf{Quando} & \textbf{Che cosa} \\ \midrule}{\bottomrule\endtabularx\par}

% --- Stili degli schemi ------------------------------------------------------------------
% Stesso stile, stesso significato, in tutto il documento.
\tikzset{
  concetto/.style={draw=mappacol,thick,fill=mappacol!12,rounded corners=4pt,align=center,text width=#1,minimum height=1cm,inner sep=5pt,font=\bfseries},
  concetto/.default=3.4cm,
  dettaglio/.style={draw=mappacol,fill=pagina,rounded corners=4pt,align=center,text width=#1,minimum height=1cm,inner sep=5pt},
  dettaglio/.default=3.4cm,
  esempio/.style={draw=libro,dashed,fill=libro!6,rounded corners=4pt,align=center,text width=#1,minimum height=1cm,inner sep=5pt},
  esempio/.default=3.4cm,
  processo/.style={draw=flussocol,fill=flussocol!8,align=center,text width=#1,minimum height=1cm,inner sep=5pt},
  processo/.default=5.5cm,
  decisione/.style={draw=ricorda,thick,fill=ricorda!10,diamond,aspect=2.2,align=center,text width=#1,inner sep=0pt},
  decisione/.default=3.4cm,
  estremo/.style={draw=flussocol,thick,fill=flussocol!18,rounded rectangle,align=center,minimum height=.9cm,inner xsep=12pt,font=\bfseries},
  freccia/.style={-{Latex[length=2.8mm,width=2.2mm]},thick,color=testo},
  collega/.style={thick,color=testo},
  etichetta/.style={font=\small,fill=pagina,inner sep=2pt,midway,align=center},
}

% --- Citazioni e testi a fronte ----------------------------------------------------------
% La lingua facoltativa deve essere dichiarata in babel.
\newenvironment{cit}[1][italian]{\begin{quote}\selectlanguage{#1}\RaggedRight}{\end{quote}}
\newcommand{\fontecit}[1]{\par{\selectlanguage{italian}#1}}
\newcommand{\afronte}[3][italian]{%
  \begin{paracol}{2}\selectlanguage{#1}\RaggedRight\sloppy #2\switchcolumn\selectlanguage{italian}\RaggedRight\sloppy #3\end{paracol}\medskip}

% --- Scritture e parole straniere --------------------------------------------------------
\newcommand{\gr}[1]{#1}
\newcommand{\en}[1]{\foreignlanguage{english}{\emph{#1}}}
% Togli il commento dopo aver dichiarato la lingua in babel:
% \newcommand{\de}[1]{\foreignlanguage{ngerman}{\emph{#1}}}
% \newcommand{\la}[1]{\foreignlanguage{latin}{\emph{#1}}}
% \newcommand{\fr}[1]{\foreignlanguage{french}{\emph{#1}}}
% \newcommand{\es}[1]{\foreignlanguage{spanish}{\emph{#1}}}
% \newcommand{\he}[1]{\foreignlanguage{hebrew}{#1}}

% --- Figure --------------------------------------------------------------------------------
% \figw{larghezza in frazione di \textwidth}{file}{didascalia}
\newcommand{\figw}[3]{\begin{figure}[H]\centering\includegraphics[width=#1\textwidth,height=.42\textheight,keepaspectratio]{#2}\caption{#3}\end{figure}}
% \figh{altezza}{file1,file2,...}{didascalia}: immagini affiancate alla stessa altezza
\newcommand{\figitem}[2]{\includegraphics[height=#1,width=.31\textwidth,keepaspectratio]{#2}\hspace{.6em}}
\newcommand{\figh}[3]{\begin{figure}[H]\centering\forcsvlist{\figitem{#1}}{#2}\unskip\caption{#3}\end{figure}}

% --- Indice --------------------------------------------------------------------------------
\newcommand{\aut}[1]{\index[indice]{#1}}
\newcommand{\idx}[1]{\index[indice]{#1}}

\title{\Huge\bfseries @@CORSO@@\\[0.6em]\Large\mdseries Appunti del corso di @@DOCENTE@@\\ @@UNIVERSITA@@}
\author{}
\date{}

\begin{document}
\frontmatter
\maketitle

\chapter*{Come leggere questi appunti}
Questi appunti contengono tutto il materiale del corso. Gli schemi e i riquadri si aggiungono al testo: non tolgono e non semplificano niente.

\textbf{Ogni capitolo ha sempre la stessa struttura.} Puoi usarla così:

\begin{enumerate}
\item Leggi \textbf{Prima di iniziare} e guarda la \textbf{Mappa del capitolo}: sai già dove stai andando.
\item Studia una sezione alla volta. Ogni sezione ha il testo e almeno uno schema.
\item Alla fine leggi \textbf{Da ricordare} e \textbf{In sintesi}.
\item Fai il \textbf{Ripasso} senza guardare, poi controlla le \textbf{Risposte}.
\end{enumerate}

\textbf{Questi segnali dicono da dove viene un'informazione.}

\begin{correzione}
Errore o incongruenza nel materiale. Il testo originale resta riportato, accanto alla correzione e al motivo.
\end{correzione}
\begin{dubbio}
Punto che le fonti del corso non permettono di risolvere: controllalo prima di usarlo all'esame.
\end{dubbio}
\begin{dallibro}
Spiegazione tratta dal libro di testo, con il numero di pagina.
\end{dallibro}
\soloappunti
Così è segnalata una sezione senza slide, ricostruita soltanto dagli appunti in classe.

\integrazione{così è segnata un'aggiunta di conoscenza generale, che non viene dal materiale del corso.}

\begin{parolesemplici}
Spiegazione aggiunta in questi appunti, dopo un passaggio difficile. Non è una frase del docente e non sostituisce il testo originale.
\end{parolesemplici}

\textbf{Alla fine del volume} trovi le schede di ripasso, il glossario e l'indice analitico. Il PDF ha i segnalibri e il testo si può far leggere da un programma di sintesi vocale.

\tableofcontents

\mainmatter
\chapter*{Mappa del corso}
\addcontentsline{toc}{chapter}{Mappa del corso}
\gdef\ultimocapitolo{Mappa del corso}\gdef\ultimasezione{}

Questa mappa mostra tutti i capitoli e i legami tra loro. Conviene tornarci ogni volta che inizi un capitolo nuovo.

\begin{schema}[il corso in una pagina]
\centering
\begin{tikzpicture}[node distance=1.1cm]
\node[concetto=6cm] (a) {Capitolo~\ref{cap:prova}: Prova tecnica};
\end{tikzpicture}
\end{schema}

\begin{center}
\begin{tabularx}{\linewidth}{@{}lY@{}}
\toprule
\textbf{Capitolo} & \textbf{Argomenti} \\
\midrule
\ref{cap:prova} & Prova del template \\
\bottomrule
\end{tabularx}
\end{center}

% Pagina di prova ed esempio d'uso di tutti gli ambienti.
% Aggiungi un campione delle scritture e dei simboli del corso, compila, guarda il risultato.
% Rileggilo come modello, poi cancella il file e la sua riga \input.
% I contenuti sono solo esempi: non vanno riportati negli appunti.
\chapter{Prova tecnica}\label{cap:prova}

\begin{primadiiniziare}
\textbf{Di cosa parla:} il capitolo mostra tutti gli strumenti del template.

\textbf{Che cosa devi ricordare:} niente, è il primo capitolo.

\textbf{Alla fine saprai:} riconoscere ogni riquadro e ogni tipo di schema.
\end{primadiiniziare}

\begin{mappa}
\centering
\begin{tikzpicture}[node distance=1.8cm and .3cm]
\node[concetto=4.4cm] (c) {Strumenti degli appunti};
\node[dettaglio,below left=of c] (s) {Schemi};
\node[dettaglio,below=of c] (t) {Testo e citazioni};
\node[dettaglio,below right=of c] (n) {Numeri e codice};
\draw[freccia] (c) -- node[etichetta] {comprende} (s);
\draw[freccia] (c) -- node[etichetta] {comprende} (t);
\draw[freccia] (c) -- node[etichetta] {comprende} (n);
\end{tikzpicture}
\end{mappa}

La mappa ha un concetto centrale, «Strumenti degli appunti», che comprende tre parti: gli schemi, il testo con le citazioni, i numeri con il codice.

\section{Schemi}

\subsection{Mappa di un concetto}

Questa è una mappa concettuale. Ogni freccia dice quale relazione lega due concetti.

\begin{schema}[idea, sue parti ed esempio]
\centering
\begin{tikzpicture}[node distance=1.8cm and 1cm]
\node[concetto=3cm] (a) {Concetto};
\node[dettaglio=3cm,right=3.8cm of a] (b) {Definizione};
\node[esempio=3cm,below=of a] (d) {Esempio del corso};
\draw[freccia] (a) -- node[etichetta] {si definisce come} (b);
\draw[freccia] (a) -- node[etichetta] {si vede in} (d);
\end{tikzpicture}
\end{schema}

Lo schema dice che un concetto ha una definizione e si vede in un esempio del corso.

\subsection{Diagramma di flusso}

\begin{flusso}[come scegliere la forma di uno schema]
\centering
\begin{tikzpicture}[node distance=1.2cm and 1.6cm]
\node[estremo] (i) {Inizio};
\node[decisione,below=of i] (d) {Ci sono passi in sequenza?};
\node[processo=4.4cm,below=of d] (f) {Usa un diagramma di flusso};
\node[processo=4.4cm,right=of d] (m) {Usa una mappa concettuale};
\node[estremo,below=of f] (e) {Fine};
\draw[freccia] (i) -- (d);
\draw[freccia] (d) -- node[etichetta] {sì} (f);
\draw[freccia] (d) -- node[etichetta] {no} (m);
\draw[freccia] (f) -- (e);
\draw[freccia] (m) |- (e);
\end{tikzpicture}
\end{flusso}

Il diagramma fa una domanda: ci sono passi in sequenza? Se la risposta è sì, si usa un diagramma di flusso. Se la risposta è no, si usa una mappa concettuale.

\section{Testo e citazioni}
% nodiagramma: sezione dimostrativa senza concetti da schematizzare

Italiano con accenti: è, à, ì, ò, ù, perché, città. Parola inglese: \en{translation}\aut{Autore, Nome}\idx{concetto}. Titolo di un'opera: \emph{Titolo dell'opera}. Termine chiave: \textbf{termine}. Nome con pronuncia: Goethe \pron{ghöte}.

Greco: \gr{παρθένος}, \gr{Ἰδοὺ ἡ παρθένος ἐν γαστρὶ ἕξει}. Parole con legature evitate: fine, fluire, affitto, ufficio. Rimando: capitolo~\ref{cap:prova}.

\begin{cit}[english]
A quotation in the language of the source stays complete, however hard it may be to read.
\fontecit{(Fonte, anno)}
\end{cit}

\begin{parolesemplici}
La citazione resta intera. Questo riquadro la spiega con parole più semplici.
\end{parolesemplici}

\afronte[english]{The source text stands in this column, in its language, across several lines.}{La traduzione sta in questa colonna, su più righe.}

\begin{nonconfondere}
\textbf{Citazione} è il testo originale. \textbf{Parafrasi} è lo stesso testo detto con altre parole.
\end{nonconfondere}

\begin{permemoria}
Un trucco per ricordare va sempre dichiarato come tale.
\end{permemoria}

\begin{correzione}Prova di correzione.\end{correzione}
\begin{dubbio}Prova di dubbio.\end{dubbio}
\begin{dallibro}Prova dal libro (p.~1).\end{dallibro}
\soloappunti
Testo. \integrazione{prova di integrazione.}

\section{Numeri, tabelle e codice}
% nodiagramma: sezione dimostrativa senza concetti da schematizzare

\subsection{Una formula}

La somma dei primi $n$ numeri è:
\begin{equation}
S = \frac{n(n+1)}{2}
\end{equation}

\begin{simboli}
$S$ & la somma cercata & nessuna \\
$n$ & l'ultimo numero da sommare & nessuna \\
\end{simboli}

\begin{passoapasso}{la somma dei primi 4 numeri}
\begin{enumerate}
\item Sostituisco $n$ con 4: $S = \frac{4 \cdot (4+1)}{2}$.
\item Calcolo la parentesi, perché ha la precedenza: $4+1 = 5$, quindi $S = \frac{4 \cdot 5}{2}$.
\item Calcolo il prodotto: $4 \cdot 5 = 20$, quindi $S = \frac{20}{2}$.
\item Divido: $S = 10$. Controllo: $1+2+3+4 = 10$.
\end{enumerate}
\end{passoapasso}

\subsection{Tabelle e date}

\begin{center}
\begin{tabularx}{\linewidth}{@{}lYY@{}}
\toprule
& \textbf{Prima colonna} & \textbf{Seconda colonna} \\
\midrule
Riga & un testo abbastanza lungo da andare a capo dentro la cella della tabella & un altro testo abbastanza lungo da andare a capo \\
\bottomrule
\end{tabularx}
\end{center}

\begin{cronologia}
1813 & Primo evento di esempio \\
1923 & Secondo evento di esempio \\
\end{cronologia}

\subsection{Codice}

\begin{lstlisting}[language=C,numbers=left]
int main(void) { /* commento con è */
    return 0;
}
\end{lstlisting}

\begin{daricordare}
\begin{itemize}
\item Ogni capitolo ha la stessa struttura.
\item Gli schemi si aggiungono al testo.
\item Ogni formula ha la spiegazione dei simboli.
\end{itemize}
\end{daricordare}

\begin{insintesi}
Il template offre schemi, riquadri e tabelle pronte. Servono a dare più vie d'accesso agli stessi contenuti.
\end{insintesi}

\begin{ripasso}
\begin{enumerate}
\item Che cosa deve avere ogni freccia di una mappa concettuale?
\item Che cosa segue sempre una citazione difficile?
\end{enumerate}
\end{ripasso}

\begin{risposte}
\begin{enumerate}
\item Un'etichetta che dice la relazione (sezione «Schemi»).
\item Un riquadro «In parole semplici» (sezione «Testo e citazioni»).
\end{enumerate}
\end{risposte}

% \input{capitoli/01-titolo}

\backmatter
% Tieni solo le schede utili al corso. Ogni riga rimanda al capitolo.
\chapter{Schede di ripasso}

Queste schede raccolgono ciò che è più difficile tenere a mente. Non sostituiscono i capitoli.

\section*{Chi ha detto cosa}

\begin{tabularx}{\linewidth}{@{}lYl@{}}
\toprule
\textbf{Chi} & \textbf{Idea} & \textbf{Capitolo} \\
\midrule
Autore & Idea principale & \ref{cap:prova} \\
\bottomrule
\end{tabularx}

\section*{Cronologia}

\begin{cronologia}
Data & Evento (capitolo~\ref{cap:prova}) \\
\end{cronologia}

\section*{Parole da non confondere}

\begin{tabularx}{\linewidth}{@{}YYl@{}}
\toprule
\textbf{Parola} & \textbf{Differenza} & \textbf{Capitolo} \\
\midrule
Citazione / parafrasi & Testo originale oppure testo riformulato & \ref{cap:prova} \\
\bottomrule
\end{tabularx}

\chapter{Glossario dei concetti}

Le voci rimandano al capitolo in cui il concetto è spiegato.

\begin{description}
\item[Concetto] Definizione breve (capitolo~\ref{cap:prova}).
\end{description}

\clearpage\gdef\ultimocapitolo{Indice analitico}\gdef\ultimasezione{}
\printindex[indice]
\end{document}
